\exo \textit{[Calculer]}

\begin{enumerate}
\item Soit $V_d=48$. Déterminer $V_f$ après une hausse de 37~\%.
\item Soit $V_f=91$. Déterminer $V_d$ après une baisse de 11~\%.
\end{enumerate}

\exo \textit{[Calculer]}

Pour chacun des coefficients multiplicateurs suivants, déterminer le taux d'évolution associé.

\setlength{\columnseprule}{0pt}
\begin{multicols}{2}
\begin{enumerate}[label=\arabic*.]
\item CM = 1,63; 
\item CM = 0,453;
\item CM = 0,9832;
\item CM = 2,45.
\end{enumerate}
\end{multicols}

\exo \textit{[Modéliser, calculer]}

Le programme ci-dessous, écrit en Python, permet de calculer le pourcentage d'évolution global, connaissant les deux pourcentages d'évolutions successives $p_1$ et $p_2$.

\medskip

\begin{center}
\begin{minipage}{0.78\linewidth}
\begin{pythoncodenumligne}
p1 = int(input("Entrer le premier pourcentage d'evolution : "))
p2 = int(input("Entrer le deuxieme pourcentage d'evolution : "))
CM1 = 1 + p1/100
CM2 = 1 + p2/100
CM = ...
p = ...
print(p)
\end{pythoncodenumligne}
\end{minipage}
\end{center}


\begin{enumerate}
\item Que faudrait-il modifier pour pouvoir saisir la valeur $p_1=+7,8$ ?  
\item Recopier l'algorithme et compléter les lignes 5 et 6.
\item Quelle est la valeur de $p$ avec $p_1=24$ et $p_2=-35$?
\end{enumerate}


